\documentclass[a4paper,12pt]{article}
\usepackage[utf8]{inputenc}
\usepackage[T1]{fontenc}
\usepackage[margin=2.5cm]{geometry}
\usepackage{ctex} % 中文支持
\usepackage{xcolor}
\usepackage{titlesec}
\usepackage{setspace}

% -----------------------------------------------------------
% 格式设置区：Level 2 标准段落 + Level 2 合规披露风格
% -----------------------------------------------------------

% 设置层级标题格式（自动编号）
\titleformat{\section}{\Large\bfseries\color{blue!40!black}}{\thesection}{1em}{}
\titleformat{\subsection}{\large\bfseries\color{gray!80!black}}{\thesubsection}{1em}{}

% 设置行间距与缩进，提升长文阅读体验
\onehalfspacing
\setlength{\parskip}{0.8em}
\setlength{\parindent}{2em}

% -----------------------------------------------------------
% 文档元数据
% -----------------------------------------------------------
\title{\textbf{2024年度环境、社会及管治（ESG）报告：可持续发展专项章节}}
\author{本公司董事会办公室}
\date{2025年3月}

\begin{document}

\maketitle

% 前言
\section*{前言：构建负责任的发展模式}

作为一家国内领先的化工行业上市公司，我们深知企业的发展不能仅看财务报表，更要看对环境、社会和治理（ESG）的综合贡献。2024年，本公司积极响应国家可持续发展战略，严格遵守各类合规监管要求，致力于构建绿色、安全、和谐的发展模式。本章节将以通俗易懂的方式，向各利益相关方详细汇报我们在过去一年的履责绩效与管理实践。

% 第一章 环境
\section{环境管理：绿色运营的合规与行动}

环境保护是公司可持续发展的底线。2024年，我们严格执行《环境保护法》及相关行业标准，通过加大资金投入和技术升级，切实降低生产运营对生态环境的影响。

\subsection{环境绩效与资源管理}

在本年度，我们建立并完善了环境数据监测体系，确保数据的真实性与可追溯性。根据实测数据，2024年公司温室气体排放总量控制在39,107.65吨CO2e，排放强度为4.58吨/百万元营收。在废弃物管理方面，我们产生一般工业固废170.48吨；对于危险废物，我们严格执行全流程管控，处置率达到100\%，确保所有危废均得到合规的安全处置。

在资源节约方面，我们坚持“取之有度，用之有节”。全年新鲜水总用量为149.55万吨，包装材料使用总量为7,677吨。为了支撑这些环保行动，公司全年环保资金总投入达到457.93万元，并依法缴纳环保税9.6万元。这些数据不仅是合规的体现，更是我们对绿色承诺的兑现。

\subsection{清洁能源与循环经济的实践}

为了从源头减少污染，我们在生产端积极推进能源结构的优化。在\textbf{核心生产基地}，我们实施了清洁能源热能替代项目。这项改造彻底淘汰了传统燃煤锅炉，转而引入清洁能源热源。这就像是将家庭取暖从烧煤球改为了使用天然气或电采暖，直接结果是该基地实现了生产环节废气的“零排放”，显著提升了厂区的环境质量。

同时，我们积极探索循环经济模式。旗下化工厂实施了冷凝水回收利用项目，将生产过程中产生的高温蒸汽冷凝水进行回收，重新作为锅炉补水。这种做法既减少了对新鲜水资源的消耗，又回收了热能，实现了节能与节水的双重效益。

\subsection{运输环节的低碳转型}

在矿山运营板块，我们重点关注移动源排放的治理。在\textbf{西北某大型露天煤矿项目}，公司推进了大规模的车辆电气化升级。我们新增了21台纯电动矿卡和36台油电增程式矿卡。这些新型车辆利用先进的能量回收技术，在下坡制动时可反向充电，大幅降低了燃油消耗。此外，我们还投入了12台无人驾驶电动矿卡。这些智能设备不仅实现了运行过程的零排放，还大幅降低了噪音污染，体现了绿色矿山的建设理念。

% 第二章 创新
\section{研发创新：驱动行业高质量发展}

科技创新是推动化工行业安全转型与效率提升的核心动力。我们坚持以技术解决实际问题，通过持续的研发投入，赋能产业升级。

\subsection{研发体系与成果转化}

2024年，公司研发投入达到4.19亿元，占营业收入的比例为4.91\%。我们组建了一支由1,064名专业人员构成的研发团队，致力于攻克行业技术难题。本年度，公司新增授权专利126项，累计拥有专利596项，这标志着我们在知识产权保护和技术积淀方面迈上了新台阶。

\subsection{关键技术突破}

我们重点推广了三项具有行业示范意义的技术成果，旨在提升本质安全水平和作业效率。

首先是“现场混装水胶炸药工艺”。这项技术达到了国际先进水平，并填补了国内空白。它的核心在于改变了炸药的形态和运输方式，将原材料分别运输至现场，在注入钻孔的瞬间混合成为炸药。这不仅消除了运输过程中的爆炸风险，还能根据岩石硬度灵活调节威力密度，实现了安全与效益的平衡。

其次是国内首创的“二氧化碳相变膨胀激发管”。针对高瓦斯矿井的开采痛点，该技术利用液态二氧化碳物理膨胀原理破岩，而非化学燃烧。整个过程无明火，极大提升了特殊环境下的作业安全性，产能已达到500万只。

最后是“智能矿山系统”。我们集成了5G、边缘计算与机器人技术，实现了5G钻机、自动验孔机器人与无人矿卡的协同作业。这一系统将工人从危险、恶劣的一线环境中解放出来，体现了“科技以人为本”的发展理念。

% 第三章 社会
\section{社会责任：以人为本与社区共建}

我们深知，企业的成功离不开员工的辛勤付出和社区的包容支持。因此，我们始终将保障员工权益和回馈社会作为履行社会责任的重点。

\subsection{员工安全与健康保障}

安全生产是化工企业的生命线。截至2024年底，公司员工总数为7,399人，劳动合同签订率保持在100\%。为了织密安全防护网，我们在安全保障上投入了大量资源，包括投入355.99万元购买安全生产责任险，以及投入892.5万元用于工伤保险。同时，我们构建了“1+14+N”的应急预案体系，确保在突发状况下能够迅速响应。

我们的安全管理水平在海外项目中得到了充分验证。在\textbf{纳米比亚}项目中，面对跨文化管理和复杂作业环境的挑战，我们建立了标准化的安全管理制度，如严格的晨会流程和“手指口述”安全确认法。凭借这些扎实的管理手段，该项目实现了连续130万小时无损工事件的优异成绩，树立了中国企业在海外安全运营的典范。

\subsection{公益慈善与应急响应}

作为央企控股企业，我们积极履行企业公民责任。2024年，公司对外捐赠及公益总投入达到1,219万元，其中1,125万元专项用于支持乡村振兴，助力改善欠发达地区的基础设施和民生水平。

在面对突发自然灾害时，我们也勇于担当。2024年\textbf{湖南某地}发生严重决堤险情，急需大量石料封堵。公司迅速响应，调集附近的矿山团队支援。依托专业的爆破技术，工程师们采用特殊的“松动爆破”方案，在极短时间内开采出8.4万吨石料，为抗洪抢险争取了宝贵时间。这一行动充分展现了公司在关键时刻的应急动员能力和社会责任感。

% 第四章 治理
\section{公司治理：合规经营与价值共享}

完善的公司治理结构是企业稳健运行的基石。作为\textbf{上级集团}的重要成员，我们致力于通过透明、高效的治理机制，保护投资者利益，维护商业道德。

\subsection{治理架构与投资者回报}

我们建立了科学的董事会结构，采用“3名股东董事+3名独立董事+3名执行董事”的平衡模式。这种结构有效保障了决策的独立性与制衡性，避免了决策风险。凭借规范的信息披露和治理实践，公司在证券交易所的信披考评中获得了“A级”评价。

在股东回报方面，我们坚持与投资者共享发展成果。2024年，公司实施了每10股派发现金红利2.05元的分配方案，共计派发2.54亿元，现金分红比例达到可分配利润的40.12\%。

\subsection{供应链合规管理}

我们高度重视供应链的廉洁建设，致力于打造公平、透明的商业环境。针对合作的1,548家供应商，我们实施了严格的准入机制，要求所有供应商必须签署《廉洁协议》。截至目前，该协议的签署率达到100\%。这一举措有效防范了商业贿赂风险，确保了采购过程的合规性，为构建健康的行业生态做出了贡献。

\end{document}